% ===== Section 2: "Do blow your fuse", anger (~18 min) =====
\topicopener{angry}{Զայրույթ}{Պայթի՛ր, կարելի է}{Ինչպես է աշխատում զայրույթը, ու ինչու է այն երբեմն օգտակար}{angry}

% Hulk
\begin{frame}{«Ինձ չես հավանի, երբ բարկացած եմ»}
  \begin{columns}[c]
    \begin{column}{0.42\textwidth}
      \centering
      \mascot[2.2cm]{angry}\par\medskip
      \begin{card}[colback=mint!40!cardwhite]
        \centering Բրյուս Բաններ՝ «Մի՛ բարկացրու ինձ։ Ինձ չես հավանի, երբ բարկացած եմ»
      \end{card}
    \end{column}
    \begin{column}{0.54\textwidth}
      \begin{itemize}
        \item<2-> Բարկանալիս Բանները դառնում է \textbf{Հալք}, որին սիրում են միլիոնները։ Ուրեմն արտահայտությունը սուտ է
        \item<3-> Բայց իսկապես, ո՞վ է սիրում բարկացած մարդուն։ Բացի «արդար զայրույթից» անարդարության դեմ
        \item<4-> Զայրույթը համարվում է վատ բան. անխոհեմ վարք, վիրավորանք, բռնություն
        \item<5-> Ինչու՞ է ուղեղն այն այդքան հեշտ արտադրում, նույնիսկ ամենաչնչին առիթով
      \end{itemize}
    \end{column}
  \end{columns}
\end{frame}

% What anger is: a violated boundary
\begin{frame}{Ի՞նչ է զայրույթը}
  \relax{\small Հուզական ու մարմնական գրգռվածություն, երբ \textbf{խախտվում է քո սահմանը}}\par\smallskip
  \begin{columns}[T]
    \begin{column}{0.31\textwidth}
      \begin{card}[colback=peach!35!cardwhite]
        \emoji[1.4em]{walking}\ {\blackfont Ֆիզիկական}\par\smallskip
        Փողոցում ինչ-որ մեկը բախվեց քեզ
      \end{card}
    \end{column}
    \begin{column}{0.31\textwidth}
      \onslide<2->{%
      \begin{card}[colback=butter!45!cardwhite]
        \emoji[1.4em]{moneybag}\ {\blackfont Ռեսուրսային}\par\smallskip
        Պարտքով փող վերցրեց ու չի վերադարձնում
      \end{card}}
    \end{column}
    \begin{column}{0.31\textwidth}
      \onslide<3->{%
      \begin{card}[colback=lavender!35!cardwhite]
        \emoji[1.4em]{judge}\ {\blackfont Բարոյական}\par\smallskip
        Քեզ համար վիրավորական կարծիք է հայտնում
      \end{card}}
    \end{column}
  \end{columns}
  \vspace{1mm}
  \onslide<4->{%
  \begin{card}[colback=brainpink!25!cardwhite]
    \textbf{Դիտմամբ} արվածը սադրանք է (provocation), ու զայրույթն ավելի ուժեղ է։ Մի բան է՝ պատահաբար թափել բաժակը, այլ բան՝ խմիչքը շպրտել դեմքին։
  \end{card}}
  \onslide<5->{\centering\small Ուղեղը տրոլներին արձագանքում էր ինտերնետից շատ առաջ \emoji{angersymbol}\par}
\end{frame}

% Recalibration theory: the stone axe
\begin{frame}{Ինչու է զայրույթը պետք եղել. քարե կացինը}
  \centering
  \begin{tikzpicture}
    \node[strip=3.3cm] (a) at (-5.1,1.3) {\emoji[1.4em]{axe}\\ Նախնին ժամերով քարե կացին է պատրաստում};
    \node[strip=3.3cm, visible on=<2->] (b) at (-1.4,1.3) {\emoji[1.4em]{monkey}\\ Մեկ ուրիշը գալիս ու վերցնում է այն};
    \node[strip=3.6cm, fill=lavender!40!cardwhite, visible on=<3->] (c1) at (2.6,2.3) {\emoji{thinking}\ Նստում ու փիլիսոփայում է սեփականության մասին};
    \node[strip=3.6cm, fill=brainpink!35!cardwhite, visible on=<4->] (c2) at (2.6,0.4) {\emoji{boxing}\ Բարկանում ու գողի ծնոտին է հարվածում};
    \node[strip=2.4cm, fill=mint!45!cardwhite, visible on=<5->, font=\footnotesize] (d) at (6.0,0.4) {Պահում է գործիքը, հարգանք ու զուգընկեր};
    \draw[arr] (a) -- (b);
    \draw[arr, visible on=<3->] (b.east) -- (c1.west);
    \draw[arr, visible on=<4->] (b.east) -- (c2.west);
    \draw[arr, visible on=<5->] (c2) -- (d);
  \end{tikzpicture}\par\smallskip
  \raggedright
  \onslide<5->{%
  \begin{itemize}
    \item \textbf{Վերակալիբրացման տեսություն} (recalibration theory). զայրույթը արագ ինքնապաշտպանություն է, երբ ինչ-որ բան ես կորցնում
    \item<6-> Էվոլյուցիոն հոգեբանությունը սիրում է պարզեցնել, ինչն ինքնին մարդկանց \textbf{բարկացնում է}
  \end{itemize}}
\end{frame}

% Threat detection
\begin{frame}{Վտանգի ահազանգը}
  \begin{columns}[T]
    \begin{column}{0.52\textwidth}
      \begin{itemize}
        \item Զայրույթը հաճախ \textbf{վտանգի} պատասխան է։ Վտանգ հայտնաբերող համակարգը. \textbf{նշաձև մարմին} (amygdala), \textbf{հիպոկամպ} ու \textbf{ջրատարի շրջակա գորշ նյութ} (periaqueductal grey)
        \item<2-> Այս հին համակարգով ենք կողմնորոշվում նաև այսօր
        \item<3-> Գործընկերը բոլորի առաջ քեզ ծաղրում է։ Վտանգված են \textbf{համբավդ ու դիրքդ}։ Արդյունքը՝ զայրույթ
      \end{itemize}
    \end{column}
    \begin{column}{0.42\textwidth}
      \onslide<4->{%
      \begin{scienceupdate}
        Գրքում այս երեքը կոչվում են «միջին ուղեղի (midbrain) մասեր»։ Իրականում միջին ուղեղում է միայն \textbf{գորշ նյութը}։ Նշաձև մարմինն ու հիպոկամպը \textbf{քունքային բլթի} խորքում են՝ լիմբիկ համակարգում։
      \end{scienceupdate}}
    \end{column}
  \end{columns}
\end{frame}

% Anger as a goal-driven emotion; how labs make people angry
\begin{frame}{Ուղեղն ինչ-որ բան է ուզում}
  \begin{columns}[c]
    \begin{column}{0.54\textwidth}
      \begin{itemize}
        \item Սկաներում բարկացած մարդկանց \textbf{ակնակապային ճակատային կեղևն} (orbitofrontal cortex) ավելի ակտիվ է (Carver \& Harmon-Jones)
        \item Այս մասը կապված է զգացմունքների կառավարման ու նպատակաուղղված վարքի հետ
      \end{itemize}
      \vspace{1mm}
      \begin{tikzpicture}[font=\footnotesize]
        \node[strip=2.1cm, visible on=<2->] (e) at (0,0) {Ինչ-որ բան պատահեց};
        \node[strip=2.1cm, visible on=<3->] (j) at (2.75,0) {Ուղեղը դժգոհ է};
        \node[strip=2.1cm, fill=brainpink!35!cardwhite, visible on=<4->] (f) at (5.5,0) {\emoji{angry} Զայրույթ $\rightarrow$ գործել};
        \draw[arr, visible on=<3->] (e) -- (j);
        \draw[arr, visible on=<4->] (j) -- (f);
      \end{tikzpicture}
    \end{column}
    \begin{column}{0.4\textwidth}
      \onslide<5->{%
      \begin{popup}{լաբ.exe}
        Հոդվածներում գրում են՝ «մասնակիցներին ներկայացվեցին զայրույթ առաջացնող խթաններ»։\par\smallskip
        Հաճախ դա նշանակում է, որ նրանց ուղղակի \textbf{վիրավորում էին}։ Կամավորներ գտնելն այդպես ավելի հեշտ է։
      \end{popup}}
    \end{column}
  \end{columns}
\end{frame}

% Myth or fact: anger lowers the stress hormone
\begin{frame}{Առասպել, թե՞ փաստ \#3}
  \mythorfact{Զայրույթը կարող է իջեցնել սթրեսի հորմոնը}{ՓԱՍՏ}{mintdark}{%
    \begin{card}[colback=mint!35!cardwhite]
      Սթրեսը վնասակար է հիմնականում \textbf{կորտիզոլի} (cortisol) պատճառով։ Օսնաբրյուկի համալսարանում Միգել Կազենի ու գործընկերների փորձերում զայրույթը \textbf{իջեցնում էր} կորտիզոլը, ու սթրեսի վնասը պակասում էր։
    \end{card}}
\end{frame}

% Two hemispheres
\begin{frame}{Սատանան ու հրեշտակը գլխիդ ներսում են}
  \centering
  % Approach vs avoidance motivation by hemisphere (book: "Do blow your fuse", Rutherford & Lindell 2011):
% the left side says "deal with it", the right side says "stay back". Brain seen from above, front at the top.
% Left bubble on click 1, right bubble on click 2.
\begin{tikzpicture}[lbl/.style={font=\small, text=plum, align=center}]
  % two hemispheres
  \filldraw[fill=amBlue!25, draw=plum, line width=1pt]
    (-0.12,1.6) .. controls (-1.4,1.7) and (-1.9,0.6) .. (-1.85,-0.2)
    .. controls (-1.8,-1.2) and (-1.0,-1.75) .. (-0.12,-1.7) -- cycle;
  \filldraw[fill=amRed!20, draw=plum, line width=1pt]
    (0.12,1.6) .. controls (1.4,1.7) and (1.9,0.6) .. (1.85,-0.2)
    .. controls (1.8,-1.2) and (1.0,-1.75) .. (0.12,-1.7) -- cycle;
  \foreach \s in {-1,1}{
    \draw[plum!60, line width=0.7pt] (\s*0.5,1.1) .. controls (\s*1.0,0.9) .. (\s*0.8,0.4);
    \draw[plum!60, line width=0.7pt] (\s*1.4,0.2) .. controls (\s*0.9,-0.1) .. (\s*1.1,-0.6);
    \draw[plum!60, line width=0.7pt] (\s*0.4,-0.6) .. controls (\s*0.8,-0.9) .. (\s*0.6,-1.3);}
  \node[font=\scriptsize, text=plum!80] at (0,1.95) {առջև};
  \node[lbl, font=\blackfont\small] at (-1.0,0.0) {Ձախ};
  \node[lbl, font=\blackfont\small] at (1.0,0.0) {Աջ};
  \node[lbl, font=\footnotesize, anchor=east] at (-2.0,-1.1) {\emoji{muscle}\ մոտեցում\\ (approach)};
  \node[lbl, font=\footnotesize, anchor=west] at (2.0,-1.1) {\emoji{fearful}\ խուսափում\\ (avoidance)};
  % bubbles
  \node[say, text width=3.3cm, callout absolute pointer={(-1.5,0.9)}] at (-4.3,1.2)
    {Ոչ, ես սա այսպես չեմ թողնի։ Պետք է լուծել};
  \node[say, text width=3.3cm, callout absolute pointer={(1.5,0.9)}, visible on=<2->] at (4.3,1.2)
    {Չէ՛, հեռու մնա, վտանգավոր է, մի վատացրու};
\end{tikzpicture}
\par\smallskip
  \onslide<3->{{\small Զայրույթը բարձրացնում է \textbf{ձախ} կիսագնդի, առաջային գոտկային կեղևի (anterior cingulate) ու ճակատային կեղևի ակտիվությունը}\par}
\end{frame}

% Anger pushes you off the diving board
\begin{frame}{Մի հրում ցատկահարթակից}
  \begin{columns}[c]
    \begin{column}{0.58\textwidth}
      \begin{itemize}
        \item Վստահ, էքստրավերտ մարդկանց մոտ հավանաբար գերակշռում է \textbf{ձախ} կողմը, նևրոտիկ կամ ինտրովերտ մարդկանց մոտ՝ \textbf{աջը}
        \item<2-> Աջը վտանգի դեմ ոչինչ չի անում, վտանգը մնում է, ու տագնապը աճում է
        \item<3-> Զայրույթն ուժեղացնում է ձախը. ինչպես երբ վարանողին \textbf{հրում են ցատկահարթակից}
        \item<4-> Կորտիզոլն իջնում է, ու տագնապը քեզ չի «սառեցնում»։ Ու զայրույթը \textbf{լավատես} է դարձնում
      \end{itemize}
    \end{column}
    \begin{column}{0.36\textwidth}
      \centering
      \bigemoji[1.3cm]{diving}\par\medskip
      \onslide<5->{%
      \begin{card}[colback=peach!35!cardwhite]
        {\blackfont Բայց}\par
        Նույն փորձերում զայրույթը \textbf{վատացնում էր} բարդ մտավոր առաջադրանքների արդյունքը։ Բարկացածն իսկապես «պարզ չի մտածում»։
      \end{card}}
    \end{column}
  \end{columns}
\end{frame}

% Negotiations, and what "let it out" means
\begin{frame}{Ուրեմն պետք է «դուրս թափե՞լ»}
  \begin{columns}[T]
    \begin{column}{0.48\textwidth}
      \begin{itemize}
        \item Բանակցություններում \textbf{տեսանելի} զայրույթը օգտակար է, նույնիսկ եթե երկու կողմն էլ բարկացած են. ավելի շատ մոտիվացիա, լավատեսություն ու անկեղծության տպավորություն
        \item<2-> Գիրքը եզրակացնում է. զայրույթը մի զսպիր, թող դուրս գա, որ սթրեսը պակասի ու գործը առաջ գնա
      \end{itemize}
    \end{column}
    \begin{column}{0.46\textwidth}
      \onslide<3->{%
      \begin{scienceupdate}
        «Դուրս թողնել» չի նշանակում \textbf{«թափել»}։ 154 հետազոտության 2024-ի վերլուծությունը ցույց տվեց, որ բարձի վրա բռունցքներով «զայրույթը թափելը» (venting) այն չի պակասեցնում, երբեմն՝ ավելացնում է։ Օգնում է հակառակը՝ \textbf{հանդարտեցնելը}. դանդաղ շնչառություն, մինչև տասը հաշվելը։
      \end{scienceupdate}}
    \end{column}
  \end{columns}
  \onslide<3->{\source{Kjærvik \& Bushman, Clinical Psychology Review, 2024}}
\end{frame}

% Count to ten: the orbitofrontal cortex as the bath overflow
\begin{frame}{Հաշվիր մինչև տասը}
  \begin{columns}[c]
    \begin{column}{0.42\textwidth}
      \begin{itemize}
        \item Զայրույթը շատ \textbf{արագ ու ուժեղ} է. դրա համար են աշխատում «հաշվիր մինչև տասը» ու «խորը շունչ քաշիր» հնարքները
        \item<2-> Ակնակապային ճակատային կեղևը զտում է զգացմունքների ազդեցությունը վարքի վրա ու արգելակում ամենավտանգավոր ազդակները
      \end{itemize}
    \end{column}
    \begin{column}{0.56\textwidth}
      \centering
      % Burnett's image for the orbitofrontal cortex: the overflow outlet on a bath with a leaky tap.
% It does not fix the leak (the anger), it stops the bath from flooding (behaviour). Labels on click 2.
\begin{tikzpicture}[lbl/.style={font=\small, text=plum, align=left}]
  % water
  \fill[amBlue!30] (-2.05,-0.95) -- (2.05,-0.95) -- (2.05,0.75) -- (-2.05,0.75) -- cycle;
  \draw[amBlue!60, line width=0.8pt, decorate, decoration={snake, amplitude=0.4mm, segment length=4mm}] (-2.05,0.75) -- (2.05,0.75);
  % tub
  \draw[plum, line width=1.4pt, rounded corners=4mm] (-2.2,1.2) -- (-2.2,-1.1) -- (2.2,-1.1) -- (2.2,1.2);
  \draw[plum, line width=1.4pt] (-2.45,1.2) -- (2.45,1.2);
  \draw[plum, line width=1.2pt] (-1.6,-1.1) -- (-1.75,-1.5) (1.6,-1.1) -- (1.75,-1.5);
  % tap with drops
  \draw[plum, line width=1.4pt, rounded corners=1mm] (-2.45,2.3) -- (-1.4,2.3) -- (-1.4,1.95);
  \fill[plum] (-1.6,2.3) rectangle (-1.2,2.45);
  \foreach \y in {1.75,1.45}{\fill[amBlue] (-1.4,\y) circle (0.06);}
  % overflow outlet
  \filldraw[fill=cardwhite, draw=plum, line width=1pt] (2.2,0.8) circle (0.13);
  \draw[-{Stealth[length=2mm]}, amBlue, line width=1pt] (2.35,0.8) .. controls (2.8,0.8) and (2.9,0.4) .. (2.9,0.0);
  % labels
  \node[lbl, anchor=south west, visible on=<2->] at (-2.45,2.5) {\textbf{զայրույթը} կաթում է};
  \node[lbl, anchor=west, visible on=<2->] at (2.95,0.3) {\textbf{ակնակապային}\\ \textbf{ճակատային}\\ \textbf{կեղև}. ավելորդը\\ դուրս է թողնում};
\end{tikzpicture}
\par\smallskip
      {\small Խնդիրը չի լուծում, բայց ջուրը չի թափվում}
    \end{column}
  \end{columns}
\end{frame}

% Rumination
\begin{frame}{Եռում ես օրերով}
  \begin{columns}[c]
    \begin{column}{0.56\textwidth}
      \begin{itemize}
        \item Զայրույթը կարող է տևել ժամեր, օրեր, նույնիսկ շաբաթներ
        \item<2-> Հիպոկամպն ու նշաձև մարմինը ստեղծում են \textbf{վառ, զգացմունքային} հիշողություններ
        \item<3-> Դեպքը մնում է հիշողության մեջ, ու մենք այն նորից ու նորից «ծամում» ենք մտքում (rumination)
        \item<4-> Այդ ժամանակ ակտիվ է \textbf{միջային նախաճակատային կեղևը} (medial prefrontal cortex). որոշումներ ու պլաններ
      \end{itemize}
    \end{column}
    \begin{column}{0.38\textwidth}
      \centering
      \begin{tikzpicture}
        \node at (0,0) {\mascot[2.6cm]{angry}};
        \node[think, visible on=<3->, font=\handfont\small] at (1.1,2.0) {...իսկ ես պիտի\\ ասեի...};
      \end{tikzpicture}
    \end{column}
  \end{columns}
\end{frame}

% Displacement
\begin{frame}{Մեկ օր, չորս փոքրիկ բան}
  \begin{columns}[c]
    \begin{column}{0.62\textwidth}
      \begin{enumerate}
        \item Ավտոմատը կերավ մանրադրամդ ու մանր չվերադարձրեց
        \item<2-> Մայրուղում մեկը վտանգավոր կտրեց ճանապարհդ
        \item<3-> Ժամը 16:56-ին ղեկավարն ասաց, որ այսօր ուշ ես մնալու
        \item<4-> Մատուցողը լատեի փոխարեն սև սուրճ բերեց
      \end{enumerate}
      \vspace{1mm}
      \onslide<2->{{\small Ոչ մեկին պատասխան չունես. եթե չես ուզում վանդալիզմ անել, վթարի ենթարկվել կամ աշխատանքից ազատվել։}\par}
    \end{column}
    \begin{column}{0.32\textwidth}
      \centering
      \scalebox{0.85}{% Displacement (book): small irritations with no outlet pile up; the next harmless target gets all of it.
% One segment per click (1-3), the waiter tips it over on click 4.
\begin{tikzpicture}
  \fill[amOrange!35, rounded corners=1.5mm] (0,0) rectangle (1.1,1.1);
  \fill[amOrange!70, visible on=<2->] (0,1.1) rectangle (1.1,2.2);
  \fill[amOrange, visible on=<3->] (0,2.2) rectangle (1.1,3.3);
  \fill[amRed, visible on=<4->] (0,3.3) rectangle (1.1,4.2);
  \draw[plum, line width=1pt, rounded corners=1.5mm] (0,0) rectangle (1.1,4.2);
  \foreach \y in {1.1,2.2,3.3}{\draw[plum!50, line width=0.6pt] (0,\y) -- (0.3,\y);}
  \node[font=\scriptsize, text=plum, rotate=90] at (-0.35,2.1) {զայրույթ};
  \node[visible on=<4->] at (0.55,4.9) {\bigemoji[1.1cm]{bomb}};
  \node[visible on=<4->, font=\blackfont\large, text=amRed, rotate=-10] at (1.9,4.75) {ԲՈՒՄ};
\end{tikzpicture}
}
    \end{column}
  \end{columns}
\end{frame}

\begin{frame}{Տեղափոխում (displacement)}
  \begin{columns}[c]
    \begin{column}{0.54\textwidth}
      \begin{itemize}
        \item Կուտակված զայրույթը ելք չունի, ու ուղեղն այն «լցնում է» առաջին հարմար թիրախի վրա, որ ճնշումը թուլանա
        \item Խեղճ մատուցողը ստանում է ամբողջ օրվա բաժինը
      \end{itemize}
      \vspace{1mm}
      \onslide<2->{%
      \begin{bookquote}
        Իհարկե, դա ամենևին չի մխիթարում նրան, ով պատահաբար բացեց կատաղության ջրարգելակը։
      \end{bookquote}}
    \end{column}
    \begin{column}{0.4\textwidth}
      \onslide<3->{%
      \begin{poll}
        Բարձրացրու ձեռքդ, եթե երբևէ \textbf{բարկացել ես մեկի վրա}, ով ընդհանրապես մեղավոր չէր։
      \end{poll}}
      \centering\onslide<3->{\bigemoji[1.2cm]{coffee}}
    \end{column}
  \end{columns}
\end{frame}

% Passive aggression
\begin{frame}{Պասիվ ագրեսիա}
  \begin{columns}[c]
    \begin{column}{0.46\textwidth}
      \begin{itemize}
        \item Կարելի է ագրեսիվ լինել առանց բռնության
        \item Սովորականից քիչ ես խոսում նրա հետ, չեզոք տոնով
        \item<2-> Բոլոր ընդհանուր ընկերներին կանչում ես, իրեն՝ ոչ
        \item<3-> Նա վստահ չէ՝ բարկացա՞ծ ես, իսկ ուղեղը \textbf{անորոշություն} չի տանում։ Պատիժ՝ առանց կանոն խախտելու
      \end{itemize}
    \end{column}
    \begin{column}{0.5\textwidth}
      \centering
      \onslide<4->{%
      \begin{meme}[6.4cm]{disaster_girl.jpg}
        \memecap{0.3}{0.88}{քո սոցիալական կյանքը}
        \memecap[text width=2.6cm]{0.76}{0.12}{«Ի՞նչ, ամեն ինչ նորմալ է»}
      \end{meme}}
    \end{column}
  \end{columns}
\end{frame}

% We read anger well
\begin{frame}{Զայրույթը հեշտ է ճանաչել}
  \begin{columns}[c]
    \begin{column}{0.5\textwidth}
      {\small Ուղեղը զայրույթը որսում է բոլոր «նուրբ» նշաններով}\par\medskip
      \begin{itemize}
        \item \emoji{people}\ մարմնի լեզու
        \item \emoji{angry}\ դեմքի արտահայտություն
        \item \emoji{loudspeaker}\ ձայնի տոն
        \item<2-> \emoji{knife}\ ժանգոտ մաչետեով քո հետևից վազելն ու գոռալը
      \end{itemize}
    \end{column}
    \begin{column}{0.44\textwidth}
      \onslide<3->{%
      \begin{card}[colback=lavender!35!cardwhite]
        Սա օգտակար է. բարկացած մարդը կարող է \textbf{վտանգավոր} լինել։\par\smallskip
        Բայց դա նաև ցույց է տալիս, որ ինչ-որ բան նրան \textbf{իսկապես վիրավորել է}։
      \end{card}}
    \end{column}
  \end{columns}
\end{frame}

% Anger, hostility, aggression
\begin{frame}{Զայրույթ, թշնամանք, ագրեսիա}
  \begin{columns}[c]
    \begin{column}{0.5\textwidth}
      \begin{tryit}
        Հարևանդ քո մեքենայի վրա հայհոյանք է գրում։ Ո՞րն է զայրույթը, որը՝ թշնամանքը, որը՝ ագրեսիան\par\smallskip
        \begin{enumerate}
          \item Զգում ես, թե ինչպես է արյունդ եռում
          \item Մտածում ես. «Սրան ջախջախելու եմ»
          \item Աղյուս ես նետում նրա պատուհանին
        \end{enumerate}
      \end{tryit}
    \end{column}
    \begin{column}{0.44\textwidth}
      \onslide<2->{%
      \begin{card}[colback=peach!35!cardwhite]
        \textbf{1. Զայրույթ} (anger). զգացմունքը, երբ քեզ սպառնում են։ Երևի բոլորի մոտ նույնն է
      \end{card}}
      \onslide<3->{%
      \begin{card}[colback=lavender!35!cardwhite]
        \textbf{2. Թշնամանք} (hostility). վնաս տալու մասին \textbf{մտքերը}
      \end{card}}
      \onslide<4->{%
      \begin{card}[colback=brainpink!30!cardwhite]
        \textbf{3. Ագրեսիա} (aggression). վնաս տալու \textbf{վարքը}։ Լինում է նաև առանց զայրույթի. ռեգբի, ֆուտբոլ
      \end{card}}
    \end{column}
  \end{columns}
\end{frame}

% Moderation: the celebrity chef
\begin{frame}{Ուրեմն բարկանա՞լ, թե՞ ոչ}
  \begin{columns}[c]
    \begin{column}{0.56\textwidth}
      \begin{itemize}
        \item Զայրույթը միշտ չէ, որ վատ է, բայց կարևորը \textbf{չափն} է
        \item<2-> Բարկացած մարդկանց խնդիրները հաճախ լուծվում են ավելի շուտ, քան քաղաքավարի խնդրողներինը
        \item<3-> Ուղեղը զայրույթը կապում է \textbf{պարգևի} հետ ու խրախուսում այն
        \item<4-> Ու ստացվում է մարդ, որ ամենափոքր անհարմարությունից պայթում է, որ իր ուզածին հասնի...
      \end{itemize}
    \end{column}
    \begin{column}{0.38\textwidth}
      \centering
      \onslide<5->{%
      \bigemoji[1.2cm]{forkknife}\ \bigemoji[1.2cm]{fire}\par\medskip
      \begin{bookquote}
        ...ու հետո, անխուսափելիորեն, դառնում է \textbf{հայտնի խոհարար}։ Լա՞վ է դա, թե՞ վատ, դու որոշիր։
      \end{bookquote}}
    \end{column}
  \end{columns}
\end{frame}
